\medskip

\emph{Cet exercice est un questionnaire à choix multiples. Aucune justification n'est attendue pour cet exercice, on en fournit quand même dans ce corrigé.}

\medskip

\begin{enumerate}
	\item \textbf{Réponse A.}

Ici, les quatre propositions peuvent se vérifier avec les longueurs des côtés du triangle ABC. On va donc les déterminer. Comme le repère est orthonormé, on a :

$\mathrm{AB} = \sqrt{(x_\mathrm{B} - x_\mathrm{A})^2+(y_\mathrm{B} - y_\mathrm{A})^2+(z_\mathrm{B} - z_\mathrm{A})^2} = \sqrt{(3 + 1)^2 + (6-2)^2 + (3 - 5)^2} = \sqrt{36}=6$

$\mathrm{BC} = \sqrt{(x_\mathrm{C} - x_\mathrm{B})^2+(y_\mathrm{C} - y_\mathrm{B})^2+(z_\mathrm{C} - z_\mathrm{B})^2} = \sqrt{(3-3)^2 + (0-6)^2 + (9-3)^2} = \sqrt{72}=6\sqrt{2}$

$\mathrm{AC} = \sqrt{(x_\mathrm{C} - x_\mathrm{A})^2+(y_\mathrm{C} - y_\mathrm{A})^2+(z_\mathrm{C} - z_\mathrm{A})^2} = \sqrt{(3 + 1)^2 + (0-2)^2 + (9 - 5)^2} = \sqrt{36}=6$

ABC est donc clairement un triangle isocèle en A, (et on peut vérifier rapidement que le carré du côté opposé à A est égal à la somme des carrés des deux autres côtés, et donc que le triangle est aussi rectangle en A).

\emph{Remarque :} si on a confiance en ses calculs, une fois qu'on a calculé AB et BC, on peut s'arrêter : on voit clairement que le triangle ne sera pas équilatéral, et donc il sera rectangle et isocèle, car c'est un point commun aux trois autres propositions et le côté [BC] étant plus long que AB, cela signifie que [BC] est l'hypoténuse, et donc que le sommet principal du triangle est A.

	\item \textbf{Réponse C.}

Ici, il faut s'armer de patience et vérifier si les coordonnées des points B, C et D vérifient les équations proposées. (On peut aussi chercher à trouver une équation du plan, mais c'est plutôt plus compliqué).

Pour la proposition \textbf{a.}, les coordonnées de D ne vérifient pas l'équation : elle est incorrecte.

Pour la proposition \textbf{b.}, les coordonnées de C et de D ne vérifient pas l'équation.

Pour la proposition \textbf{d.}, aucunes des coordonnées des points B, C et D ne vérifient pas l'équation.

Par contre, pour l'équation de la proposition \textbf{c.}, on a :

$4 x_\mathrm{B}+y_\mathrm{B}+z_\mathrm{B}-21=4\times 3 + 6 + 3 - 21 = 12 + 9 - 21 =0$ donc B appartient à ce plan.

$4 x_\mathrm{C}+y_\mathrm{C}+z_\mathrm{C}-21=4\times 3 + 0 + 9 - 21 = 12 + 9 -21 =0$ donc C appartient à ce plan.

$4 x_\mathrm{D}+y_\mathrm{D}+z_\mathrm{D}-21=4\times 8 - 3 - 8 - 21 = 32 - 11 - 21 = 0$ donc D appartient à ce plan.

Ainsi, l'équation de la proposition \textbf{c.} décrit un plan qui contient les points B, C et D : c'est le plan (BCD).

	\item \textbf{Réponse B.}

Ici, on peut procéder par tests successifs : vérifier parmi les coordonnées proposées quelles sont celles qui vérifient l'équation du plan (ABC) puis pour celles qui vérifient l'équation du plan (c'est-à-dire quand le point donné est bien sur le plan) si le vecteur $\vect{\mathrm{DH}}$ obtenu est bien un vecteur normal au plan.

Dans ce corrigé, on va appliquer la démarche qui sert à trouver les coordonnées de H.

Comme le plan (ABC) a pour équation cartésienne $x-2 y-2 z+15=0$, cela signifie que le vecteur $\vect{u} \begin{pmatrix}1\\-2\\-2\end{pmatrix}$ est un vecteur normal au plan. Comme H est le projeté orthogonal de D sur (ABC), H est l'intersection du plan (ABC) avec la droite $d$, passant par D et dirigée par $\vect{u}$.

Du coup, la droite $d$ admet une représentation paramétrique qui est : $\left\{\begin{array}{l} x=8+k\\y=-3-2k\\z=-8-2k \end{array}\right. \quad k \in \R$

Si on appelle $M_k$ le point de paramètre $k$ sur la droite $d$, on va chercher pour quel paramètre $M_k$ est sur le plan (ABC) :

$\begin{aligned} M_k \in (\mathrm{ABC}) &\iff (8+k) -2(-3-2k) - 2--8-2k)+15=0\\
&\iff 8+k+6+4k+16+4k+15 = 0\\
&\iff 9k + 45 = 0\\
&\iff k = -5\\ \end{aligned}$

H est donc le point $M_{-5}$ qui a comme coordonnées : $(8+(-5)~;~-3-2\times (-5)~;~-8-2\times (-5))$, soit $(3 ~;~ 7 ~;~ 2)$.
	\item \textbf{Réponse D.}

La représentation paramétrique de $\Delta$ donne un vecteur directeur de $\Delta$ qui est $\vect{v}$, de coordonnées $\begin{pmatrix}1\\-1\\3\end{pmatrix}$ alors que le vecteur $\vect{\mathrm{BC}}$ a pour coordonnées $\vect{\mathrm{BC}}\begin{pmatrix}0\\-6\\6\end{pmatrix}$. Ces vecteurs sont clairement non colinéaires, donc on peut affirmer que les droites se sont ni confondues ni strictement parallèles.

Il faut donc déterminer si elles ont un point commun pour trancher entre sécantes et non coplanaires.

La droite (BC) passe par B et est dirigée notamment par $\dfrac{1}{6}\vect{\mathrm{BC}}$, elle admet pour représentation paramétrique : $\left\{\begin{array}{l} x=3\\y=6-s\\z=3+s \end{array}\right. \quad k \in \R$

Le point de paramètre $t$ sur $\Delta$ est confondu avec celui de paramètre $s$ sur (BC) si et seulement si $s$ et $t$ sont les solutions du système $(S)$ suivant :

\setlength{\columnseprule}{0.8pt}
\begin{multicols}{2}
$(S) : \left\{\begin{array}{l} 3 = 5+t\\6-s=3-t\\3+s=-1+3t \end{array}\right.
\iff \left\{\begin{array}{l} t = -2\\6-s=5\\3+s=-7 \end{array}\right.	$

$(S) \iff  \left\{\begin{array}{l} t = -2\\s=1\\s=-10 \end{array}\right.$
	\end{multicols}

Ce système n'a pas de solution, donc les droites n'ont aucun point commun et ne sont pas parallèles : elles sont donc non coplanaires.

\item \textbf{Réponse B.}

Le plan $\mathcal{P}$ d'équation cartésienne $2 x-y+2 z-6=0$ admet $\vect{n_1} \begin{pmatrix}	2\\-1\\2\end{pmatrix}$ comme vecteur normal.

(ABC) d'équation cartésienne $x-2 y-2 z+15=0$ admet $\vect{n_2} \begin{pmatrix}1\\-2\\-2\end{pmatrix}$ comme vecteur normal.

Ces vecteurs sont clairement non colinéaires, donc les plans ne sont pas parallèles. On peut même calculer le produit scalaire $\vect{n_1}\cdot\vect{n_2}$ pour constater qu'ils sont orthogonaux, et donc que les plans sont perpendiculaires. Cela élimine la proposition \textbf{a.}

Reste à trancher entre les trois autres propositions.

On pourrait résoudre le système formé par les deux équations de plan pour trouver une représentation paramétrique de leur intersection : ce serait lourd en calcul. Ici, il est plus simple de tester l'appartenance des points A, B et C au plan $\mathcal{P}$. En remplaçant $x$, $y$ et $z$ par les coordonnées des points dans l'équation de $\mathcal{P}$, on trouve que les points A, et B appartiennent à $\mathcal{P}$, mais pas C, donc les plans sont sécants et leur intersection est la droite (AB).
	\end{enumerate}

\bigskip
